\scanpage{41}
Soit $A$ une VA sur $k$ alg. clos,
d'où module $\mathbb T_\ell(A)\deleted{\unsure{\otimes}}$ sur la pro-alg. de Lie $\deleted{\unsure{\mathfrak g_K}}\unsure{\mathfrak g|\mathscr G}(k)=\mathfrak g$.
Soit $K$ une clôture alg. de $\mathbb Q_\ell$, posant
\[
\mathbb T_{K}(A)=\mathbb T_\ell(A)
\]
de sorte que $\mathbb T_K(A)$ est un module sur
$\mathfrak g\otimes_{\mathbb Q_\ell}K=\mathfrak g_K$. Supposons

a) La \deleted{\unsure{hyp}} conj de Tate vérifiée pour $H^2(A,\mathbb Q_\ell(1))$
\[
H^2(A\deleted{)},\mathbb Q_\ell(1))^{\mathrm{alg}}=H^2(A,\mathbb Q_\ell(1))^{\mathfrak g}
\]

b) $\mathfrak g_K$ opérant sur $\deleted{\unsure{\mathfrak g}}\deleted{\text{est}}\mathbb T_K(A)$ est ``pseudoalgébrique''
[i.e. les valeurs propres des frobenius \deleted{sont} \unsure{de|des}
norme $q$ sont $\varepsilon q^{-1/2}$, $\varepsilon$ une racine de l'unité \ldots]

% slanted
\margin{\unsure{Par|Pour} Tits, le premier cas est un cas particulier du second}
\uline{ou} $\mathbb T_K(A)$ se décompose en deux morceaux
$V_1\times V_2$, tels que chacun soit propre pour ``presque
tous'' les Frobénius, avec des valeurs propres
différentes pour presque tous les Frobénius.

[Ces conditions sont remplies si $A$ est une \deleted{produit}\inserted{puissance $C^n$}
\deleted{de}\inserted{d'une} courbe elliptique \inserted{$C$ définissable} sur un corps fini $\mathbb F_q$,
comme on voit en notant que a) \deleted{\inserted{\unsure{résulte}} et b)
sont stables par passage à une puissance
cartésienne, et pour une $C$, a) est}
a été vérifié par Mumford (on est ramené
% archive sheet number: 40
